\section{为什么意识研究应当置于此处}

原始框架只把行为状态当作实践层面的对象来处理。
对于早期概念稿而言，这已经足够，但它会使全书与当前生物学研究之间保持过远距离。
如果意识通达、清醒水平、注意力以及大尺度协调本身也是类似状态的生物学现象，
那么关于状态持续与状态转变的语言，就不应停留在日常生产力层面。

因此，本章明确遵循手稿的证据政策。
本章的\textbf{强实证核心}，是文献地图中分配给本章的官方索引来源：
\emph{Dynamical structure-function correlations provide robust and generalizable signatures of consciousness in humans}
（Communications Biology, 2024）、
\emph{Cortical connectivity, local dynamics and stability correlates of global conscious states}
（Communications Biology, 2025），以及
\emph{Falling asleep follows a predictable bifurcation dynamic}
（Nature Neuroscience, 2025）。
\textbf{中等}层是解释性的：像 Stanislas Dehaene 的 \emph{Consciousness and the Brain}
与 Christof Koch 的 \emph{Consciousness} 这样的经典著作，有助于说明这些发现为何对本书重要，
但它们并不用于替代实证核心。
\textbf{推测性}的意识物理学方案仍然不属于本章范围，而应放在附录~\ref{app:frontier}。